\subsection{19.62: the derived carpet is closed}
\label{cand:19.62}

Let \(\Phi\) be a reduced crystallographic root system and \(K\) a
commutative ring with identity. A carpet \(A=(A_r)_{r\in\Phi}\) is a
family of additive subgroups of \(K\) satisfying all Chevalley conditions
\[
 c_{ij,rs}a^ib^j\in A_{ir+js}\quad
 (a\in A_r,\ b\in A_s,\ i,j>0,\ ir+js\in\Phi).
\]
Here \(r,s\) are distinct and nonopposite, and \(c_{ij,rs}\) is the
absolute value of the corresponding integral commutator constant.
Signs do not affect additive subgroup membership. Write
\(\partial A=B\) for the family additively generated by these
parameters at each destination root. Then \(B\subseteq A\) and \(B\)
is itself a carpet: a rule applied to elements of \(B\) is also a rule
applied to elements of \(A\), so its value belongs to the appropriate
subgroup of \(B\).

The candidate proves the stronger inclusion
\begin{equation}
 A_p^2B_{-p}\subseteq B_p \qquad(p\in\Phi),
 \label{eq:derived-square}
\end{equation}
where \(A_p^2=\{a^2:a\in A_p\}\). Consequently
\(B_p^2B_{-p}\subseteq B_p\). Nuzhin's already published
square-closedness criterion then says that \(B\) is closed:
\[
 \langle x_r(B_r):r\in\Phi\rangle\cap x_p(K)=x_p(B_p).
\]
The criterion, and hence the answer to \kn{19.63}, predates the
experiment \cite[Theorem~2]{nuzhin-2026}. The ledger's grouping of
19.62 with 19.63 must not be read as a new solution of the latter.

Substantial cases of \kn{19.62} were also known. Nuzhin's 2023
discussion records the answer for type \(A_\ell\); the earlier
triple-product inclusion quoted there, together with its Theorem~1,
also gives the simply laced types \(D_\ell,E_\ell\).
Its Corollary~2 treats \(B_\ell,C_\ell,F_4\) when
\(\gcd(\operatorname{char}K,2)=1\), and \(G_2\) when
\(\gcd(\operatorname{char}K,6)=1\) \cite[Section~3]{nuzhin-2023}.
The integer derivations below apply without these characteristic
restrictions; the already covered cases are not new answers.

\paragraph{Reduction to finite integer derivations.}
Fix \(x\in A_p\). It suffices to prove
\[
 c x^2u^iv^j\in B_p
 \quad\text{whenever}\quad
 ir+js=-p,\quad u\in A_r,\quad v\in A_s,
 \tag{\(\ast\)}
\]
since multiplication by \(x^2\) distributes over every additive
expression for an element of \(B_{-p}\).
The intersection of \(\Phi\) with the plane spanned by \(r,s\) is an
irreducible rank-two root system, of type \(A_2,B_2\), or \(G_2\).
Every root string between its roots stays in that plane.
Thus the restricted carpet has precisely the required absolute
Chevalley constants, and any derived parameter in the restriction
belongs to the ambient \(B\). No commutation of restriction with
the entire operation \(\partial\) is assumed.

The following finite proof calculus establishes \((\ast)\).
Initial nodes are independent variables with specified membership
in \(A_r\). A Chevalley rule applied to two earlier nodes gives
\(c f^i g^j\in B_{ir+js}\subseteq A_{ir+js}\).
If two applications give integer multiples \(m_1w,m_2w\) of the
same monomial at the same root, their integer linear combinations
give \(\gcd(m_1,m_2)w\in B\). Each final coefficient must divide
the requested coefficient. These operations use no division in \(K\)
and remain valid in characteristics two and three.

There are 6, 20, and 78 distinct requests in types \(A_2,B_2,G_2\),
respectively, after coincident requests are combined by taking the
gcd of their coefficients. Short certificates contain 655 nodes in
total, at most nine per request, including initial nodes. Relabeling
the carpet by Weyl transformations reduces the display to 19 requests,
one destination root of each length. This is transport of a universal
implication; it does not assert that the particular carpet is Weyl
invariant. The ancillary checker verifies the absolute-rule invariance,
the coverage by those representatives, and all 104 requests directly.
The 19 representative derivations are printed in
Appendix~\ref{app:carpet-tables}.

For illustration, take \(G_2\) with short simple root \(a\), long
simple root \(b\), and \(p=2a+b\). For
\(x\in A_p,u\in A_a,v\in A_{-3a-b}\), the rules give
\[
 3xu\in B_{3a+b},\quad xv\in B_{-a},\quad
 x^2v\in B_{a+b}.
\]
Applying the rules again gives \(3x^2uv\in B_p\) from the first
two expressions and \(2x^2uv\in B_p\) from the third and \(u\).
Their difference gives \(x^2uv\in B_p\). This illustrates why
retaining both small integer coefficients can prove an assertion
that an argument dividing by a structure constant would miss.

\paragraph{A direct proof of the known closedness criterion.}
We retain the following alternative argument as part of the
experiment's mathematical record. Assume \(A_r^2A_{-r}\subseteq A_r\)
for every root and work in the integral Chevalley Lie algebra tensored
with \(K\). Put
\[
 M=\sum_r A_re_r+
       \sum_r\langle A_rA_{-r}\rangle_{\Z}h_r.
\]
Root coordinates in this additive group are independent; in particular
\(M\cap Ke_r=A_re_r\). The group \(M\) is a Lie subring.
Root--root brackets use the \((1,1)\) carpet condition or, for
opposite roots, the displayed Cartan summands. For a Cartan generator
\(uvh_r\) and a root \(s\ne\pm r\), Jacobi expresses their bracket
as two successive root brackets, each allowed by the carpet rules.
For \(s=r\) or \(-r\), polarization supplies the required coefficient:
\[
 2uvw=(u+w)^2v-u^2v-w^2v\in A_r
 \quad (u,w\in A_r,\ v\in A_{-r}).
\]
The Cartan part commutes with itself. This also shows that \(M\)
is exactly the Lie subring generated by the root summands.

Every \(x_s(t)\), \(t\in A_s\), preserves \(M\).
On \(ue_r\) with \(r\ne\pm s\), its root-string terms have
coefficients \(C_{i1,sr}t^iu\), covered by the carpet rules.
It fixes \(e_s\), while
\[
 ue_{-s}\longmapsto ue_{-s}+tu h_s-t^2u e_s
\]
uses the square assumption. The inverse generator gives equality.
If a root \(p\) admits a root \(q\) with
\(\langle p,q^\vee\rangle=\pm1\), all generating root elements
send \(h_q\) into \(h_q+M\). Thus \(x_p(t)\) in the carpet group
forces \(\pm te_p\in M\), and hence \(t\in A_p\).
Such a \(q\) exists in simply laced systems of rank at least two
by adjacency; in \(B_n\), \(n\ge3\), use \(e_1+e_2\) for a short
root \(e_1\), and \(e_2+e_3\) for a long root \(e_1+e_2\).
The same choices cover \(F_4\), and the same-length hexagons cover
\(G_2\).

For \(C_n\), including \(B_2=C_2\) and rank one, use natural
symplectic matrices indexed by \(1,\ldots,n,-1,\ldots,-n\).
The root vectors are
\[
 E_{ij}-E_{-j,-i},\qquad E_{i,-j}+E_{j,-i},\qquad E_{i,-i},
\]
and their negative-root counterparts. Together with
\(H_i=E_{ii}-E_{-i,-i}\) they form a free Chevalley basis over \(K\).
Every root vector \(N\) satisfies \(N^2=0\).
Choose a diagonal matrix \(D\) with entries \(d_i\in\{0,1\}\)
at \(i\) and \(1-d_i\) at \(-i\). Then
\[
 [D,e_r]=\lambda_re_r,\qquad
 \lambda_{e_i-e_j}=d_i-d_j,\quad
 \lambda_{e_i+e_j}=d_i+d_j-1,\quad
 \lambda_{2e_i}=2d_i-1.
\]
For any prescribed \(p\), the \(d_i\) can make \(\lambda_p=\pm1\).
Since \(NDN=0\),
\[
 (I+tN)D(I-tN)=D-t\lambda_rN.
\]
Each initial root generator sends \(D\) into \(D+M\), and invariance
of \(M\) extends this to every word. The formula detects \(t\in A_p\)
for a new root element \(I+te_p\).
This also works in every standard Chevalley isogeny form:
an equality of root words gives equality of actions on the
abstract Chevalley Lie algebra. The quotient of their natural
symplectic lifts centralizes that algebra. Commuting with
\(E_{i,-i}\) and \(E_{-i,i}\) forces a matrix to be scalar on each
two-coordinate block and zero between blocks; the short-root
matrices force the scalars to agree. Such a quotient is therefore
scalar over any commutative ring and commutes with \(D\).
The same detection argument applies. This proves the criterion
without dividing by two.

The retained proof and audits are
\artifact{research/19.62-rank-two-proof.md},
\artifact{research/19.62-closedness-criterion.md}, and
\artifact{results/19.62-rank-two-packet.json}.
The portable ancillary files include the proof calculus, target
coverage checks, and integer certificates.
